% Lettre classique, sans questions ni sections, police Helvetica comme le CV.
\documentclass[12pt,a4paper]{article}
\usepackage[a4paper,top=1.9cm,bottom=1.9cm,left=2.4cm,right=2.4cm]{geometry}
\usepackage[utf8]{inputenc}
\usepackage[T1]{fontenc}
\usepackage[scaled=1.0]{helvet}
\renewcommand{\familydefault}{\sfdefault}
\usepackage[hidelinks]{hyperref}
\usepackage{xcolor}
\definecolor{muted}{HTML}{555F67}
\setlength{\parindent}{0pt}
\setlength{\parskip}{0.65em}
\linespread{1.25}
\emergencystretch=2em
\pagestyle{plain}
\clubpenalty=10000
\widowpenalty=10000
\brokenpenalty=10000
\begin{document}
\thispagestyle{empty}
\textbf{Stéphane KPOVIESSI}\\
{\small Cergy, Île-de-France\\
+33 7 48 59 47 13\\
\href{mailto:oastephanekpoviessi@gmail.com}{oastephanekpoviessi@gmail.com}}

\vspace{0.4em}
Comité d'admission\\
Faculté des sciences de l'administration\\
Université Laval

\vspace{0.3em}
\textbf{Objet : Candidature au MBA en finance -- Automne 2027}

Madame, Monsieur,

Mon intérêt pour la finance s'est précisé en travaillant sur les données des entreprises. Savoir mesurer une activité ne suffit pas à déterminer quel projet financer, quelle valeur attendre d'un investissement ou quel risque accepter. Ingénieur diplômé de JUNIA ISEN en Big Data et Data Science, je souhaite approfondir ces questions au MBA en finance de l'Université Laval.

Chez Allianz France, j'ai contribué à la migration d'un datamart d'assurance et à la fiabilisation des indicateurs associés. Comprendre les règles métier et contrôler les données m'a appris à examiner ce qu'un chiffre représente avant de l'utiliser. Cette expérience a nourri mon intérêt pour les mécanismes financiers de l'assurance. Chez Lyne \& Frères, les analyses commerciales et opérationnelles m'ont rapproché des préoccupations d'une direction : suivre la performance, comprendre les écarts et disposer d'informations fiables pour agir. Je souhaite désormais acquérir les méthodes permettant d'en apprécier les conséquences financières.

Mon parcours scientifique constitue un appui pour cet apprentissage. J'ai suivi un seul cursus d'ingénieur de cinq années, de 2020 à 2025, au Bénin, au Maroc puis en France. Les deux années préparatoires m'ont donné des bases solides en mathématiques ; ma spécialisation m'a ensuite formé au traitement des données et à l'apprentissage automatique. Mon expérience en vision par ordinateur chez Sylvagreg a renforcé ma rigueur expérimentale. Je souhaite mettre ces acquis au service de l'analyse financière, tout en apprenant les concepts économiques et comptables nécessaires pour interpréter les résultats d'un modèle.

Je choisis le MBA en finance de l'Université Laval pour sa formation spécialisée en finance d'entreprise, en analyse financière et en gestion des risques, inscrite dans une perspective de gestion. L'enseignement d'Économétrie financière et analyse de données offre un lien particulièrement intéressant avec mon parcours quantitatif. Parmi les cours proposés, Gestion financière de projets rejoint mon souhait d'apprendre à évaluer un investissement et à comparer des options de financement.

À court terme, je souhaite exercer comme analyste financier ou analyste de la performance, notamment dans les services ou l'assurance. Je souhaite contribuer à l'évaluation des projets et aux choix d'investissement, avec une double compétence en finance et en données. À moyen terme, j'ambitionne de prendre davantage de responsabilités dans le pilotage financier et les orientations de développement d'une organisation.

Depuis septembre 2026, je développe également Seshat, un projet de plateforme d'IA générative et d'automatisation. Définir une offre et choisir les solutions qui la soutiennent m'amène à arbitrer des priorités de développement. Je souhaite apprendre à en examiner la viabilité économique, les besoins de financement et les perspectives de croissance. À plus long terme, ces compétences doivent aussi me permettre de contribuer au développement de Lyne \& Frères, l'entreprise familiale active au Bénin, en évaluant ses projets avec une compréhension plus complète de leurs enjeux financiers.

Mes responsabilités en communication puis à la vice-présidence de JUNIA ALUMNI m'ont initié au travail collectif et à la coordination. Je souhaite apporter aux échanges à Laval cette expérience, ma perspective d'ingénieur et les questions concrètes rencontrées dans une grande entreprise comme dans une structure familiale. Je vous remercie de l'attention portée à ma candidature et me tiens à votre disposition pour tout renseignement complémentaire.

Veuillez agréer, Madame, Monsieur, l'expression de mes salutations distinguées.

\vspace{0.5em}
\textbf{Stéphane KPOVIESSI}
\end{document}
